\fancyhead[L]{\small KLÜBER ÜBER DIE STERNWARTE ZU MANNHEIM.}
\hypertarget{native-541-lower}{}\label{native-541-lower}
\begin{center}\small Göttingische gelehrte Anzeigen. Stück 192 und 193. S. 1913..1916. 2. Dezember 1811.\\[3pt]\rule{.72\linewidth}{.3pt}\end{center}

\emph{Die Sternwarte zu Mannheim, beschrieben von ihrem Kurator, dem Staats- und Kabinettsrat} \textsc{Klüber}. Mit einer Abbildung der Sternwarte im Steindruck. 1811. 62 Seiten in Kleinfolio. In Kommission bei \textsc{Gottlieb Braun}. Mannheim und Heidelberg.

{\rightskip=0pt plus 1em
Die Mannheimer Sternwarte ist gegenwärtig unstreitig von allen Sternwarten in Deutschland am reichsten mit Instrumenten ausgestattet. Jeder, dem die Astronomie etwas bedeutet, nimmt lebhaften und herzlichen Anteil an ihr. Ein ausführlicher Bericht über ihre Entstehung, ihre Entwicklung und ihren gegenwärtigen Zustand kann daher nur willkommen sein. Doppelt erfreulich ist der Bericht aus der Feder ihres Kurators, dem der Ruf und das Wohl des Instituts so sehr am Herzen liegen und der dies durch sein Handeln bereits so deutlich gezeigt hat.\par}

Das Fundament für das durchweg massive Gebäude wurde 1772 unter dem Kurfürsten von der Pfalz, \textsc{Carl Theodor}, nach dem Plan des bekannten Astonomen P. \textsc{Mayer} gelegt. Um den Bau vollkommen dauerhaft und schön auszuführen, scheute man keine Kosten. Sie beliefen sich auf über 70000 Gulden. Das wird niemanden wundern, wenn man die beträchtliche Höhe von 111 Fuß bedenkt. Sie galt damals als notwendig, muss aber bei der heutigen Art zu beobachten als wesentlicher Fehler der Anlage gelten, da sie das Beobachten ebenso erschwert, wie sie seine Zuverlässigkeit beeinträchtigt. Die Sternwarte wurde mit Instrumenten der besten englischen Instrumentenbauer ausgestattet. Ein hervorragender achtfüßiger Mauerquadrant von \textsc{Bird} wurde schon 1775 aufgestellt. Im Jahr 1778 erhielt die Sternwarte einen zwölffüßigen Zenitsektor von \textsc{Sisson} und einen Regulator von \textsc{Arnold}. Später lieferte \textsc{Ramsden} das sechsfüßige Meridianfernrohr (für $145\tfrac{1}{2}$ Guineen). Ein dreifüßiger Multiplikationskreis von \textsc{Reichenbach} in München wurde erwartet und ist privaten Nachrichten zufolge inzwischen tatsächlich eingetroffen. Der erste Astronom, \textsc{Christian Mayer}, starb 1783 (sein durch Spezialtafeln für Aberration und Nutation bekannter\SourceLeaf{542} Gehilfe, \textsc{Joh. Metzger}, starb 1780). Von 1784 bis 1786 leitete \textsc{Carl König} die Sternwarte, in den Jahren 1786 und 1787 \textsc{Joh. Nepomuk Fischer}. Dessen vorgesehener Nachfolger, \textsc{Peter Ungeschick}, starb vor Antritt der Stelle. Seit 1788 ist schließlich \textsc{Roger Barry} Astronom an der Sternwarte (eine Zeit lang gemeinsam mit \textsc{Henry}). Der Krieg war für die Arbeiten sehr ungünstig. Ende des Jahres 1794 mussten alle Instrumente abgenommen und in Kisten verpackt werden, in denen sie 6 Jahre blieben Mehrmals wurde die Sternwarte beim Bombardement Mannheims sogar von Haubitzen getroffen, ohne jedoch wesentlich beschädigt zu werden. Der Astronom war persönlichen Misshandlungen und Verfolgungen ausgesetzt, denen er sich nur durch Flucht entziehen konnte. Nach dem Frieden von Lunéville, als Mannheim an das Haus Baden fiel, wurden zwar mehrere wissenschaftliche Sammlungen und Kunstsammlungen von Mannheim nach München gebracht, doch blieben die kostbaren astronomischen Instrumente Mannheim erhalten und wurden 1801 wieder aufgestellt. Erst später wurden die nötigsten Mittel für den Unterhalt der Sternwarte und die Besoldung des Astronomen bereitgestellt.

Im Jahr 1808 wurde die Sternwarte der unmittelbaren Fürsorge des Staatsrats \textsc{Klüber} anvertraut, der alles tat, um ihr wieder aufzuhelfen. Die nötigen Mittel für den Unterhalt, die Verschönerung und die Möblierung des Gebäudes wurden bewilligt, damit es auch gelegentlich auswärtige Astronomen beherbergen konnte. Bewilligt wurden ferner Mittel für Beleuchtung, Heizung, Korrespondenz, eine Bibliothek und andere kleinere Bedürfnisse. Auf der Nordseite wurde ein schöner Obelisk für die Mire méridienne aufgestellt, und ein zweiter wird gegenwärtig auf der Südseite errichtet. Die Regierung hat die Mittel für den Druck des großen Sternkatalogs angewiesen, an dem \textsc{Barry} seit vielen Jahren arbeitet. Ein Multiplikationskreis von \textsc{Reichenbach} wurde, wie oben erwähnt, erst kürzlich angeschafft. Der Astronom selbst erhielt eine beträchtliche Gehaltserhöhung und zugleich einen bezahlten Aufwärter. Gegenwärtig wird auch ein Assistent für den Astronomen in Betracht gezogen.

Die badische Regierung und der verdiente Kurator der Sternwarte verdienen gewiss den Dank aller Freunde dieser Wissenschaft für die großzügige Unterstützung einer schönen Einrichtung, die der praktischen Astronomie gewidmet ist. Von ihrer tatkräftigen und kundigen Nutzung darf man zu Recht reiche Früchte erwarten.

Der übrige Teil der vorliegenden Schrift gibt einen Überblick über die bisherigen Ortsbestimmungen der Sternwarte, deren Länge von Paris auf $24^{\mathrm m}31^{\mathrm s}8$ und deren Breite auf $49^\circ\,29'\,16''$ festgesetzt wird. Er enthält außerdem ein Verzeichnis aller vorhandenen Instrumente, biografische Nachrichten über die bisherigen Astronomen der Sternwarte und ein vollständiges Verzeichnis ihrer Schriften.
